\thispagestyle{empty}
\vspace*{2.0cm}

\begin{center}
{\Large\bfseries\color{titlecolor} LỜI MỞ ĐẦU\par}
\end{center}

\vspace{1.05cm}

Quyển 36 ở bản trước có hai vấn đề dễ thấy: dàn trang bị vụn vì tách từng câu quá nhiều, và phần thơ chưa đủ tự nhiên để tạo tiếng cười hay độ mượt khi đọc liền mạch. Vì vậy, bản này được làm lại theo một nguyên tắc đơn giản hơn: \textbf{thơ phải đọc xuôi trước, đẹp trang sau, rồi mới đến phần song ngữ}.\par

Nội dung vẫn giữ tinh thần học trò, nhưng giọng được kéo về gần sân trường hơn: có ghẹo nhau, có chối tội, có lười học, có sợ kiểm tra, có rung động rất nhỏ, và có cả những màn nghịch dại khiến thầy cô vừa bực vừa buồn cười. Thay vì cố gồng thành những câu quá kiểu cách, các bài được viết theo nhịp tự nhiên, sáng nghĩa, vui miệng và dễ nhớ hơn.\par

Về bố cục, sách được chia thành 12 chương nhỏ, mỗi chương 5 bài, để người đọc đi từng lát cắt cảm xúc: thích thầm, bạn cùng bàn, giờ học, trò nghịch, mùa thi, phao thi, giám thị gọi tên, tan học đợi nhau. Mỗi bài được đóng trong một khối gọn, phần tiếng Việt đứng trung tâm, còn phần tiếng Anh chỉ là một dòng gloss ngắn để giữ tính song ngữ mà không làm bài thơ bị vỡ nhịp.\par

Hy vọng lần chỉnh này khiến quyển 36 đọc vui hơn, mềm hơn và hợp với không khí học trò hơn: mở ra là thấy có tiếng lớp học, tiếng trống trường, tiếng ghẹo nhau, và cả chút đỏ mặt rất quen của tuổi áo trắng.

\vfill
\begin{flushright}
\textit{Tháng 3 năm 2026}
\end{flushright}

\clearpage
